\chapter{Command Combinations and Playbooks}
\label{ch:play}
A command is a word; a playbook is a sentence. Each recipe below is complete, ordered, and says what you
should see after every step.

\section{Playbook 1: cold start to first dashboard}
\begin{lstlisting}[style=console]
$ sudo kbctl ssh add-key ~/.ssh/id_ed25519.pub     # once per operator
$ sudo kbctl up
$ kbctl status                                      # all four active
$ kbctl dashboard open                              # browser appears
$ kbctl stats
events/sec:       112.40
active processes: 238
\end{lstlisting}
If \texttt{stats} errors with \texttt{dial unix /run/kb/kba.sock: no such file}, \cmd{kbd} is not running:
\texttt{systemctl status kbd}.

\section{Playbook 2: a remote operator takes over}
On the server (once): \texttt{sudo kbctl ssh add-key} for their key, \texttt{sudo kbctl ssh up}.
On their laptop:
\begin{lstlisting}[style=console]
$ ssh -t -p 2222 operator@kb-host        # lands directly in the TUI, nothing else
\end{lstlisting}
They need no account on the server and get no shell. Revoke by deleting their line from
\file{/etc/kb/authorized_keys}; no restart is needed.

\section{Playbook 3: incident, from detection to containment}
\begin{figure}[ht]
\centering
\begin{tikzpicture}[x=1cm,y=1cm,scale=0.9,transform shape]
  \node[box,minimum width=2.7cm] (d) at (0,0) {1. DETECT\\[-1pt]\scriptsize dashboard / TUI alert};
  \node[box,minimum width=2.7cm] (i) at (3.7,0) {2. INVESTIGATE\\[-1pt]\scriptsize \texttt{pid}, MCP, audit};
  \node[dark,minimum width=2.7cm] (c) at (7.4,0) {3. CONTAIN\\[-1pt]\scriptsize \texttt{process isolate}};
  \node[box,minimum width=2.7cm] (v) at (11.1,0) {4. VERIFY\\[-1pt]\scriptsize \texttt{audit verify}};
  \draw[flow] (d)--(i); \draw[flow] (i)--(c); \draw[flow] (c)--(v);
  \draw[flowd] (v.south) -- ++(0,-0.9) -| node[lbl,pos=0.25,below]{released? \texttt{--level} lower or \texttt{none}} (c.south);
\end{tikzpicture}
\caption{The response loop. Containment is reversible until you choose \texttt{TERMINATE}.}
\end{figure}
\resetsteps
\step{Detect.} A process enters \texttt{BORDERLANDS}. The dashboard raises a \texttt{BORDERLANDS\_ENTRY} alert.
\step{Investigate.} In the TUI Query Console: \texttt{pid 4242}, \texttt{zone borderlands}.
\step{Contain, mildest effective level first.}
\begin{lstlisting}[style=console]
$ sudo kbctl process isolate --pid 4242 --level SECCOMP \
      --reason "beaconing to unknown ASN, ticket 1187"
pid=4242 contained at level SECCOMP
\end{lstlisting}
Escalate to \texttt{NAMESPACE} only if the behaviour continues; reserve \texttt{TERMINATE} for confirmed hostile code.
\step{Verify the record.}
\begin{lstlisting}[style=console]
$ kbctl audit verify
chain intact — 1432 entries verified     # counts will differ
\end{lstlisting}

\section{Playbook 4: change policy without restarting}
\label{sec:policy-change}
\begin{lstlisting}[style=console]
$ sudo $EDITOR /etc/kb/policy.yaml
$ sudo kbctl policy reload
\end{lstlisting}
\cmd{kbd} re-reads the file it was \emph{started with} (\file{/etc/kb/policy.yaml} in the packaged unit) and
logs how many process policies and extra sensitive paths it loaded.

\section{Playbook 5: protect a workload with signed CWP rules}
\label{sec:cwp}
Protected-workload rules can require a signature so nobody can weaken them by editing a file.
\resetsteps
\step{Create a signing key pair somewhere safe (not on the monitored host if you can avoid it).}
\begin{lstlisting}[style=console]
$ kbctl workload keygen --out-dir ~/kb-keys --name cwp-signer
Wrote private key (keep secret): /home/you/kb-keys/cwp-signer.key
Wrote public key (for kbd's --workloads-pubkey): /home/you/kb-keys/cwp-signer.pub
\end{lstlisting}
\step{Edit and sign. The file is rewritten \emph{in place}.}
\begin{lstlisting}[style=console]
$ sudo $EDITOR /etc/kb/workloads.yaml
$ sudo kbctl workload sign /etc/kb/workloads.yaml \
      --key ~/kb-keys/cwp-signer.key --signed-by "alice@soc" --version 2
Signed /etc/kb/workloads.yaml (signed_by="alice@soc" version=2)
\end{lstlisting}
\step{Tell \cmd{kbd} to require the signature.} Install the public key and add a drop-in (Section~\ref{sec:dropin}):
\begin{lstlisting}[style=console]
$ sudo install -m644 ~/kb-keys/cwp-signer.pub /etc/kb/workloads.pub
$ sudo systemctl edit kbd
[Service]
ExecStart=
ExecStart=/usr/bin/kbd --db /var/lib/kbd/state.db --policy /etc/kb/policy.yaml \
    --rules /etc/kb/rules.yaml --workloads /etc/kb/workloads.yaml \
    --workloads-pubkey /etc/kb/workloads.pub
$ sudo systemctl restart kbd
\end{lstlisting}
\step{Push live to every connected sensor.}
\begin{lstlisting}[style=console]
$ sudo kbctl workload reload
\end{lstlisting}
Without \texttt{--workloads-pubkey}, CWP stays unsigned and opt-in; an unsigned file is accepted.

\section{Playbook 6: prove the audit log is intact}
\label{sec:audit}
\begin{lstlisting}[style=console]
$ kbctl audit verify
chain intact — 1432 entries verified           # counts will differ; exit 0
$ kbctl audit export --out /tmp/audit-$(date +%F).json
wrote 1432 entries to /tmp/audit-2026-10-27.json
\end{lstlisting}
Each entry hashes the previous one (SHA-256). A broken chain prints
\texttt{CHAIN BROKEN — N entries verified before break}. Export before you investigate a suspected break so
you hold a copy of the evidence.

\section{Playbook 7: scripting and health gates}
\begin{lstlisting}[style=console]
# wait until the SSH console really answers, then continue
$ until kbctl ssh status >/dev/null; do sleep 1; done

# fail a CI step if the dashboard is not serving
$ kbctl dashboard status || { echo "dashboard down"; exit 1; }

# watch the stack
$ watch -n2 kbctl status

# stop everything for maintenance, start it again
$ sudo kbctl down && sudo kbctl up

# one-liner health report
$ kbctl status; kbctl stats; kbctl audit verify
\end{lstlisting}
Exit code 3 from the two \texttt{status} commands means ``not up''; code 1 means the command itself failed.

\section{Playbook 8: a quick demo for an audience}
\begin{lstlisting}[style=console]
$ sudo kbctl up
$ kbctl dashboard open                       # screen 1: live dashboard
$ kbctl ssh connect                          # screen 2: the same data in a terminal
# in the TUI: Tab to Processes, press / to filter, c to contain a demo process
$ kbctl audit verify                         # screen 3: the audit proof
\end{lstlisting}
Remember the TUI's \texttt{[DEMO DATA]} banner when \cmd{kbd} is down: do not present fictional rows as live.

\section{Combination cheat sheet}
\begin{center}\small
\begin{tabularx}{\textwidth}{@{}L{5.6cm}Y@{}}
\toprule
\textbf{Goal} & \textbf{Command}\\
\midrule
Everything up and verified & \texttt{sudo kbctl up \&\& kbctl status}\\
Web only & \texttt{kbctl dashboard open}\\
TUI only, on this box & \texttt{kbctl tui}\\
TUI over SSH, loopback & \texttt{kbctl ssh connect}\\
TUI over SSH, remote & \texttt{kbctl ssh connect HOST -i KEY}\\
Add an operator & \texttt{sudo kbctl ssh add-key KEY.pub}\\
Contain, then prove it & \texttt{kbctl process isolate ... \&\& kbctl audit verify}\\
Policy edit, apply & \texttt{sudo \$EDITOR /etc/kb/policy.yaml \&\& sudo kbctl policy reload}\\
Evidence copy & \texttt{kbctl audit export --out FILE}\\
Stop for maintenance & \texttt{sudo kbctl down}\\
Stop one surface permanently & \texttt{kbctl dashboard down} / \texttt{sudo kbctl ssh down}\\
\bottomrule
\end{tabularx}
\end{center}
